% ID: SEL-D0025-OA-E02
% ESTADO: validada. S026, cotejada con examen y criterios.
\textbf{a)} El dominio es $[0,\infty)$. Los dos tramos son
continuos y derivables en sus interiores. En $t=2$ el valor
y ambos límites valen $4$. Las pendientes laterales son $4$
y $-4$, así que no hay derivada en $2$. En $0$ hay continuidad
y derivada por la derecha, igual a $0$.

\textbf{b)} En $(0,2)$, $f'(t)=2t>0$ y en $(2,\infty)$,
$f'(t)=-4/(t-1)^2<0$. Crece hasta $2$ y decrece desde allí.
Como $4/(t-1)\to0$ para $t\to+\infty$, la asíntota horizontal
es $y=0$. El valor $t=1$ pertenece al tramo polinómico y
no es una asíntota vertical.
\begin{center}
\setlength{\unitlength}{9mm}
\begin{picture}(10,6)(-1,-0.7)
\qbezier(0,0)(4,0)(8,0)\put(8.2,0){$t$}
\qbezier(0,0)(0,2.3)(0,4.6)\put(0,4.7){$f$}
\qbezier(0,0)(1,0)(2,4)
\qbezier(2,4)(2.333,2.667)(3,2)
\qbezier(3,2)(3.667,1.333)(5,1)
\qbezier(5,1)(6,0.75)(8,0.5714)
\put(2.1,4.2){$(2,4)$}
\put(3.1,2.05){$(3,2)$}
\put(5.1,1.08){$(5,1)$}
\end{picture}
\end{center}
\textbf{c)} El cambio de creciente a decreciente sitúa el
máximo en $\boxed{2}$ segundos. La contracción es de
$\boxed{4}$ décimas de milímetro, es decir, $0.4$ mm.
